\begin{statement}{E15}{Proved}{Tangent cohomology and Kuranishi germs of positive combs}
For each positive-comb manifold $X_m$ constructed in
\cite{Viaclovsky2026Combs}, in particular for every $m\geq12$ in
its stated construction, one has
\[
 h^\bullet(X_m,T_{X_m})=(2m+1,\,8m+1,\,6m+1,\,0).
\]
Its Kuranishi germ is smooth and reduced of dimension $8m+1$.
A locally miniversal family is supplied by all sufficiently nearby
admissible period sections on the root curve, with the affine
attachment classes fixed. Its members have algebraic dimension one
and torus fibrations with the same three critical values,
multiplicities $3,4$, and hexagonal fibre type.
\end{statement}
\begin{proof}
Let $\mathfrak B=\mathbb P(3,4)$, with coarse projection
$\rho:\mathfrak B\to B$. Write $p_1,p_2$ for the finite-isotropy
points and $p_0$ for the hexagonal value. On the root curve let
$\mathcal H$ be the extended cohomology bundle, $F$ the Hodge
subbundle, $G$ its quotient, and $V$ its dual:
\[
 0\longrightarrow F\longrightarrow\mathcal H
 \longrightarrow G\longrightarrow0,\qquad V=F^*.
\]
The finite-cover comparison identifies
$G=R^1\bar f_*\mathcal O_{X_m}$ by invariant pushforward of the
ordinary torus Hodge sequence. At the hexagon the two deck
characters of the infinite toric quilt map to the two generators
of $H^1(\mathcal O_W)$ in its normalization complex, establishing
the same comparison there. The bundle $\bar f^*V$ is the vertical
logarithmic tangent bundle: its translations are dual to the
holomorphic one-forms on smooth fibres, and at the hexagon the
dual logarithmic coordinate fields extend in every toric chart.
Cup product identifies $R^2\bar f_*\mathcal O$ with
$D:=\bigwedge^2G$. On finite covers this is the ordinary torus
exterior algebra; at the hexagon it is the exterior algebra of
the two cycles in the normalization complex.

We specify finite characters by the orders of their invariant
generators. At $p_1$ the nontrivial Hodge line has invariant
generator $w$ times an equivariant frame, while the nontrivial
line of each of $V$ and $G$ has generator $w^2$ times its frame.
At $p_2$ the corresponding orders are $3$ for $F$ and $1$ for
$V,G$. Each bundle also has one trivial-character line. These
characters persist because comb smoothing fixes the limiting
periods. The tautological quotient on each tooth is
$\mathcal O(1)^{\oplus2}$, so each tooth adds $2$ to the coarse
degree of $G$. Therefore
\[
 \deg\rho_*G=2m-1,\qquad D=\mathcal O_{\mathfrak B}(24m-1).
\]
The determinant of $\mathcal H$ has degree zero on the root
curve: its determinant local system is trivial, and its cusp
residue has trace zero. Thus
\[
 \det F=\mathcal O(1-24m),\quad
 \det V=\mathcal O(24m-1),\quad
 \deg\rho_*V=2m-1.
\]
Indeed the finite fractional correction for $V$ and $G$ is
$2/3+1/4=11/12$; subtracting it from $(24m-1)/12$ gives
$2m-1$. For $F$ the correction $1/3+3/4=13/12$ gives
$\deg\rho_*F=-2m-1$.

The coarse bundles $V_c=\rho_*V$ and $G_c=\rho_*G$ have
nonnegative splitting. To prove this, evaluate the unchanged
invariant real homology class $\delta$ on $F$. Its resulting
holomorphic section of $V$ never vanishes on a smooth fibre:
a real class annihilating a Hodge plane also annihilates its
conjugate. The same argument holds on the finite good-reduction
covers. At the hexagon, $\delta$ is a nonzero combination of
angular cycles and the logarithmic forms have identity angular
period matrix. Its evaluation is nonzero there too, and gives
a nowhere-zero section of $V_c$. The invariant real cohomology
class $\gamma$ likewise gives a nowhere-zero section of $G_c$;
at the cusp it belongs to $\ker N$, whose projection to the
Hodge quotient is an isomorphism in the transverse graph chart.
At each finite point these nonzero invariant central values give
nonzero values of the coarse sections. A split rank-two bundle
of positive degree on $\mathbb P^1$ with a negative summand
cannot have a nowhere-zero section: every section lies in its
other summand, which has positive degree and hence has zeros.
Since $2m-1>0$, both splittings are nonnegative. In particular
\begin{equation}\label{eq:E15-positive}
 h^0(V_c)=h^0(G_c)=2m+1,\qquad H^1(V_c)=H^1(G_c)=0.
\end{equation}

Set $A=\rho_*(V\otimes G)$ and $L=\mathcal O_{\mathfrak B}(-5)$.
Multiplication of invariant sections gives exact comparisons
\begin{equation}\label{eq:E15-products}
\begin{split}
 0&\longrightarrow V_c\otimes G_c\longrightarrow A
       \longrightarrow\mathbb C_{p_1}\longrightarrow0,\\
 0&\longrightarrow V_c(2m-1)\longrightarrow\rho_*(V\otimes D)
       \longrightarrow\mathbb C_{p_1}\longrightarrow0,\\
 0&\longrightarrow G_c(-1)\longrightarrow\rho_*(L\otimes G)
       \longrightarrow\mathbb C_{p_1}\longrightarrow0.
\end{split}
\end{equation}
Here is the local verification. At $p_1$, multiplying two
nontrivial invariant generators of $V,G$ gives $w^4$, whereas
their tensor character has generator $w$; multiplication is
therefore one factor $t=w^3$. At $p_2$ the product is $w^2$
and the tensor generator also has order two, giving no defect.
Products involving a trivial character have no defect. The
line $D$ has the nontrivial character of $G$ and coarse degree
$2m-1$, proving the second comparison. The invariant-generator
orders of $L$ are $1,1$, and $\rho_*L=\mathcal O_B(-1)$.
Its product with the nontrivial generator of $G$ has the same
single defect at $p_1$, proving the third comparison.

All summands on the left in \eqref{eq:E15-products} have zero
$H^1$, by \eqref{eq:E15-positive}; for the last sequence their
degrees are at least $-1$. Consequently
\[
\begin{array}{c|c|c}
\text{bundle on }B&h^0&h^1\\ \hline
 A&8m+1&0\\
 \rho_*(V\otimes D)&6m&0\\
 \rho_*(L\otimes G)&2m&0
\end{array}
\]
by Riemann--Roch: their degrees are $8m-3$, $6m-2$, $2m-2$
and their ranks $4,2,2$. Also
$\rho_*(L\otimes D)=\mathcal O_B(2m-1)$.
Finite-group invariant pushforward is exact, so these calculations
give the cohomology on the root curve.

The logarithmic tangent sequence is
\[
 0\longrightarrow\bar f^*V\longrightarrow T_{X_m}(-\log W)
 \longrightarrow\bar f^*L\longrightarrow0.
\]
Projection and Leray yield
\[
 h^\bullet(\bar f^*V)=(2m+1,8m+1,6m,0),\qquad
 h^\bullet(\bar f^*L)=(0,2m,2m,0).
\]
Its only possible connecting map between nonzero groups is
$H^0(L\otimes G)\to H^0(V\otimes D)$. This is the map of
sections of the fibrewise cup-product boundary in E5. The cusp
has rank two, so its sheaf map is generically injective. A map
from a locally free sheaf on a smooth curve with zero generic
kernel is injective; hence this map on sections is injective.
It follows that
\[
 h^\bullet(T_{X_m}(-\log W))=(2m+1,8m+1,6m,0).
\]
The branch-normal quotient is $\eta_*K_{dP_6}$, since the toric
three-form trivializes the canonical bundle near the fibre.
Its cohomology is $(0,0,1)$, and its exact sequence proves the
asserted tangent-cohomology tuple.

We construct a family with this full tangent space. All nearby
root-period sections have normal bundle $V\otimes G$, tangent
space $H^0(A)$, and obstruction space $H^1(A)=0$. This section
germ can be realized using an ordinary coarse Grassmannian bundle
and finite analytic conditions. Split the equivariant bundle as
$F\oplus G$ on a finite-cover chart and write a nearby plane as
a graph matrix. Dividing the invariant-generator factors shows
that a holomorphic lift requires one entry to vanish at $p_1$
and two at $p_2$. At $p_1$ the entry is from the nontrivial Hodge
line to the trivial quotient line. At $p_2$ it is that entry
and the entry between the two nontrivial lines. These are linear
conditions on the values of the local graph matrix. Their normal
sheaf is $A$. The Cech gluing obstruction across a small Artin
extension lies in $H^1(A)=0$; each local graph chart and its
stated conditions is smooth. The finite-dimensional Douady
section germ with these conditions is therefore smooth of
dimension $8m+1$. This description does not require a global
manifold cover of the root curve.

These sections produce compact threefolds even though their
limiting periods are allowed to vary. At a finite point the
equivariant lift is a small deformation on the compact
good-reduction cover. Admissibility is open there. The affine
action has the same real formula $x\mapsto Tx+v/n$; it is
holomorphic for the new equivariant periods and stays free
because its real action is unchanged. At the hexagon the
transverse graph chart is open, and its periods are
\[
 Z(q)=\frac{\log q}{2\pi i}I+H(q)
\]
with $H$ holomorphic and uniformly bounded on a smaller parameter
polydisc. The translations with unit factors
$\exp(2\pi iH_{ij})$ preserve the same toric fan. On logarithmic
radii their action is
$r\mapsto r+(I+O(1/\log|q|))\lambda$; it is uniformly proper
and free for small $q$. A fixed finite collection of toric unit
polydiscs gives compact representatives. Thus the filling is
proper with local equation $q=x_1x_2x_3$, and its total space
is smooth over the parameter. Varying $H(0)$ varies gluing units
and preserves these estimates and local equations. On the
compact complement of the critical discs admissibility is open.
The annular gluings extend holomorphically with their fixed
integral classes: the lattice--Lie--torus exact sequence and
Cartan vanishing on an annular product extend the original
translations. This constructs a proper holomorphic family with
all $8m+1$ period directions and the specified central member.

Its Kodaira--Spencer map is injective. A first-order trivialization
would make its map variation a section of
$H^0(X_m,f^*T_B)=H^0(B,T_B)$. The three critical values remain
constant by the equations $w^3,w^4,x_1x_2x_3$, so that base
vector field vanishes at three points and is zero. The
trivialization would then be fibre preserving. On a marked
smooth torus the period derivative is the corresponding
constant Dolbeault tensor in
$\operatorname{Lie}(F)\otimes H^1(F,\mathcal O_F)$.
A nonzero constant tensor cannot be a Dolbeault coboundary,
since derivatives of periodic coefficients integrate to zero.
Thus a period section giving a trivial total-space deformation
vanishes on every regular fibre, and hence identically. Its
Kodaira--Spencer map is an isomorphism onto the already computed
$H^1(T_{X_m})$. The inverse-function argument of E13 then makes
the Kuranishi ideal zero, proving smoothness, reducedness and
miniversality. The topology theorem of the comb construction
identifies these manifolds with $S^6$; E1 and E6 give algebraic
dimension one throughout their small deformation germs.
\end{proof}

\paragraph{Comparison with the constructed parameter space.}
The germ in \cite[Lemma labelled \texttt{paramspace} and
Proposition labelled \texttt{effective}]{Viaclovsky2026Combs}
fixes the values of the coarse period section at all three
critical points. Its normal sheaf is $N(-p_0-p_1-p_2)$, where
$N=(\rho_*F)^*\otimes G_c$. The actual inclusion of that germ
into our root-period germ has the differential induced by
\[
0\longrightarrow N(-p_0-p_1-p_2)\longrightarrow A
\longrightarrow
 \mathbb C_{p_0}^{4}\oplus\mathbb C_{p_1}^{3}
                     \oplus\mathbb C_{p_2}^{2}
\longrightarrow0.
\]
The local graph conditions above prove this exact sequence:
the old sheaf sets all four entries to zero at each point,
whereas $A$ imposes zero, one, and two conditions, respectively.
For the published combs its left term has zero $H^1$, as proved
in the cited smoothing construction. Its tangent dimension is
$8m-8$, and the quotient contributes exactly nine directions.
This explains the relation of the published lower bound to the
dimension $8m+1$ proved here by an explicit morphism of parameter
germs. The source asks for the full Kuranishi germ in
Section~10(1). Searches on 15 November 2026 for its title together
with ``Kuranishi'', ``smooth'', and ``$8m+1$'', and for the
two-parameter and Liu--Rollenske titles with ``tangent cohomology'',
located the supplied construction papers but no earlier complete
computations E13--E15. This records the literature comparison;
we claim no priority.


% E19 from part_E19_topology.tex

